\begin{table}[h]
\centering\small
\begin{tabular}{@{}lrr@{}}
\toprule
configuration & $F_x$ & ripple \\
\midrule
fixed body, bulk velocity $+U$  & $+5.66034$ & $\pm0.00000$ \\
body at $+U$, bulk velocity $0$ & $-5.69040$ & $\pm0.15625$ \\
\midrule
relative difference & \multicolumn{2}{r}{$0.53\%$} \\
\bottomrule
\end{tabular}
\caption{Galilean invariance. $U=0.5$, $\nu=1$, $n_x=81$, $\dt=0.025$, no slip; each case run to
$t=30$ and averaged over the last $400$ steps. The moving case carries a ripple of $\pm2.8\%$ at
the frequency $U/h$ at which the body slides from one cut pattern to the next --- the
characteristic signature of a non-body-fitted cut-cell method, and the same effect that makes the
cylinder drag a poor instrument for measuring the order of convergence. The means agree to well
inside it.}
\label{tab:galilean}
\end{table}
